% Source: conv_net_evolutions.tex
\begin{notechapter}
\chapter{The Evolution of CNNs}\label{ch:cnn-evolution}
\chaptersubtitle{From LeNet to R-CNN}

\section{LeNet (1998)}
\begin{itemize}
  \item \textbf{Key Contribution}: Introduced by Yann LeCun, LeNet was 
    one of the first successful CNNs, designed for handwritten digit 
    recognition (e.g., MNIST dataset).
  \item \textbf{Architecture}:
    \begin{itemize}
      \item Convolutional layers followed by subsampling (pooling) 
        layers.
      \item Fully connected layers for classification.
    \end{itemize}
  \item \textbf{Impact}: Demonstrated the effectiveness of CNNs for 
    image recognition, though limited by computational resources and 
    small datasets at the time.
\end{itemize}

\section{AlexNet (2012)}
\begin{itemize}
  \item \textbf{Key Contribution}: Developed by Alex Krizhevsky, Ilya 
    Sutskever, and Geoffrey Hinton, AlexNet revolutionized deep 
    learning by winning the ILSVRC with a significant margin.
  \item \textbf{Architecture}:
    \begin{itemize}
      \item Deeper network with 5 convolutional layers and 3 fully 
        connected layers.
      \item Introduced \textbf{ReLU activation} for faster training.
      \item Used \textbf{dropout} to reduce overfitting.
      \item Employed \textbf{data augmentation} and \textbf{GPU 
        acceleration} for training.
    \end{itemize}
  \item \textbf{Impact}: Popularized deep learning and CNNs, showcasing 
    their potential for large-scale image classification.
\end{itemize}

\section{VGG (2014)}
\begin{itemize}
  \item \textbf{Key Contribution}: Developed by the Visual Geometry 
    Group at Oxford, VGG emphasized the importance of depth in CNNs 
    using small 3x3 filters.
  \item \textbf{Architecture}:
    \begin{itemize}
      \item Stacked multiple 3x3 convolutional layers to increase 
        depth.
      \item Used 2x2 max-pooling layers for downsampling.
      \item Popular variants: VGG-16 and VGG-19 (16 and 19 layers).
    \end{itemize}
  \item \textbf{Impact}: Formalized CNN design with a simple, modular 
    structure, demonstrating that depth improves performance. It set a 
    benchmark for CNNs and influenced future architectures.
\end{itemize}

\section{Inception (GoogLeNet, 2014)}
\begin{itemize}
  \item \textbf{Key Contribution}: Introduced by Google, the Inception 
    network (GoogLeNet) won the ILSVRC 2014 with a focus on efficiency 
    and reducing computational cost.
  \item \textbf{Architecture}:
    \begin{itemize}
      \item Introduced the \textbf{Inception module}, which uses 
        multiple filter sizes (1x1, 3x3, 5x5) in parallel to capture 
        features at different scales.
      \item Used \textbf{1x1 convolutions} for dimensionality reduction 
        and computational efficiency.
      \item Implemented \textbf{auxiliary classifiers} to combat 
        vanishing gradients.
    \end{itemize}
  \item \textbf{Impact}: Demonstrated the importance of network 
    efficiency and multi-scale feature extraction. Later versions 
    (e.g., Inception-v2) incorporated \textbf{Batch Normalization}
    \footnote{Batch Normalization is a technique that normalizes the 
    inputs of each neuron in a layer across a mini-batch of data. It 
    stabilizes and accelerates training by reducing internal covariate 
    shift and allowing higher learning rates. Covariate shift refers 
    to the change in the distribution of inputs to a model during 
    training. Batch Normalization mitigates this by ensuring that the 
    distribution of activations remains stable, allowing for faster 
    convergence and improved generalization.} to stabilize training.
\end{itemize}

\section{ResNet (2015)}
\begin{itemize}
  \item \textbf{Key Contribution}: Developed by Kaiming He et al., 
    ResNet (Residual Network) won the ILSVRC 2015 and introduced the 
    concept of \textbf{residual learning} to enable training of very 
    deep networks.
  \item \textbf{Architecture}:
    \begin{itemize}
      \item Introduced \textbf{skip connections} (or residual 
        connections) that allow gradients to flow directly through the 
        network, mitigating the vanishing gradient problem.
      \item Enabled the training of networks with hundreds of layers 
        (e.g., ResNet-152).
      \item Used \textbf{Batch Normalization} to stabilize and 
        accelerate training.
    \end{itemize}
  \item \textbf{Residual Function}:
    \begin{itemize}
      \item In a residual layer, the output \( \mathbf{y} \) is:
        \[
        \mathbf{y} = f(\mathbf{x}, \mathbf{W}) + \mathbf{x},
        \]
        where:
        \begin{itemize}
          \item \( \mathbf{x} \) is the input to the layer,
          \item \( f(\mathbf{x}, \mathbf{W}) \) is the residual 
            function, which learns the difference between the input and 
            the desired output,
          \item \( \mathbf{W} \) represents the weights of the layer.
        \end{itemize}
      \item The residual function \( f(\mathbf{x}, \mathbf{W}) \) 
        typically consists of multiple layers, such as:
        \[
        f(\mathbf{x}, \mathbf{W}) = \text{Conv2}(\text{ReLU}(
        \text{Conv1}(\mathbf{x}, \mathbf{W}_1)), \mathbf{W}_2),
        \]
        where \( \text{Conv1} \) and \( \text{Conv2} \) are 
        convolutional layers, and \( \text{ReLU} \) is the activation 
        function.
    \end{itemize}
  \item \textbf{Why Residual Connections Help}:
    \begin{itemize}
      \item \textbf{Gradient Flow}:
        \begin{itemize}
          \item During backpropagation, the gradient of the loss \( 
            \mathcal{L} \) with respect to \( \mathbf{x} \) is:
            \[
            \frac{\partial \mathcal{L}}{\partial \mathbf{x}} = 
            \frac{\partial \mathcal{L}}{\partial \mathbf{y}} \cdot 
            \left( \frac{\partial f(\mathbf{x}, \mathbf{W})}{
            \partial \mathbf{x}} + \mathbf{I} \right),
            \]
            where \( \mathbf{I} \) is the identity matrix.
          \item The term \( \mathbf{I} \) ensures that gradients can 
            flow directly through the skip connection, preventing them 
            from vanishing in very deep networks.
        \end{itemize}
      \item \textbf{Degradation Problem}:
        \begin{itemize}
          \item In very deep networks, adding more layers can sometimes 
            degrade performance. Residual connections address this by 
            allowing the network to learn small perturbations 
            (residuals) rather than full transformations.
        \end{itemize}
      \item \textbf{Easier Optimization}:
        \begin{itemize}
          \item By learning residuals, the network can focus on 
            refining the input rather than learning the entire 
            transformation from scratch, making optimization easier.
        \end{itemize}
    \end{itemize}
  \item \textbf{Impact}: Revolutionized deep learning by making it 
    possible to train extremely deep networks, achieving 
    state-of-the-art performance on many tasks.
\end{itemize}

\section{R-CNN Family (2014--2017)}
\begin{itemize}
  \item \textbf{Key Contribution}: A family of models for object 
    detection and instance segmentation, combining region proposals 
    with CNNs.
  \item \textbf{Deep Features}:
    \begin{itemize}
      \item R-CNNs rely on \textbf{deep features}\footnote{Deep 
        features are high-level representations of an image learned by 
        a CNN. They capture semantic information about objects, such as 
        edges, textures, shapes, and object parts, and are used for 
        tasks like classification, bounding box regression, and 
        instance segmentation.} extracted from CNNs to represent 
        region proposals.
      \item These deep features capture high-level semantic information 
        about objects, enabling accurate classification, bounding box 
        regression, and instance segmentation.
    \end{itemize}
  \item \textbf{R-CNN (2014)}:
    \begin{itemize}
      \item Uses an external algorithm (e.g., Selective Search) to 
        generate region proposals.
      \item Extracts deep features from each region proposal using a 
        pre-trained CNN (e.g., AlexNet).
      \item Classifies objects using an SVM and refines bounding boxes 
        using regression.
      \item \textbf{Limitation}: Computationally expensive due to 
        independent processing of region proposals.
    \end{itemize}
  \item \textbf{Fast R-CNN (2015)}:
    \begin{itemize}
      \item Processes the entire image through a CNN to generate a 
        feature map.
      \item Uses RoI pooling to extract deep features for each region 
        proposal from the feature map.
      \item Enables end-to-end training, improving efficiency.
    \end{itemize}
  \item \textbf{Faster R-CNN (2016)}:
    \begin{itemize}
      \item Replaces external region proposal algorithms with a 
        \textbf{Region Proposal Network (RPN)}.
      \item The RPN generates region proposals directly from the 
        feature map, using deep features to predict \textbf{objectness 
        scores}\footnote{The objectness score is a measure of how 
        likely a region proposal contains an object of interest (as 
        opposed to background). It is predicted by the RPN and used to 
        rank region proposals. This is distinct from \textbf{IoU 
        (Intersection over Union)}, which measures the overlap between 
        a predicted bounding box and a ground truth box.} and refine 
        bounding boxes.
      \item Trains the RPN end-to-end with the rest of the network, 
        further improving speed and accuracy.
    \end{itemize}
  \item \textbf{Mask R-CNN (2017)}:
    \begin{itemize}
      \item Extends Faster R-CNN to perform \textbf{instance 
        segmentation}\footnote{Instance segmentation is a computer 
        vision task that combines object detection and semantic 
        segmentation. It identifies and delineates each distinct object 
        of interest in an image at the pixel level, assigning a unique 
        label to each instance of an object.} by adding a mask 
        prediction branch.
      \item Uses RoI Align to extract deep features more accurately, 
        preserving spatial information.
      \item The mask branch uses deep features to predict pixel-level 
        masks for each object.
    \end{itemize}
  \item \textbf{Impact}: Revolutionized object detection and instance 
    segmentation, setting new benchmarks for accuracy and efficiency.
\end{itemize}

\section{Summary of Evolution}
\begin{itemize}
  \item \textbf{LeNet}: Proved CNNs could work for image recognition.
  \item \textbf{AlexNet}: Scaled CNNs to deeper architectures and large 
    datasets, using GPUs and ReLU.
  \item \textbf{VGG}: Formalized CNN design with small 3x3 filters and 
    emphasized depth, setting a benchmark for CNNs.
  \item \textbf{Inception}: Improved efficiency and multi-scale feature 
    extraction with the Inception module, later adopting BatchNorm.
  \item \textbf{ResNet}: Enabled training of very deep networks through 
    residual connections and BatchNorm, achieving groundbreaking 
    performance.
  \item \textbf{R-CNN Family}: Revolutionized object detection and 
    instance segmentation by combining region proposals with CNNs and 
    leveraging deep features.
\end{itemize}
\end{notechapter}
